\clearpage\begingroup\SourceLeaf{396b}\ArithmeticHead{65}{R}{DIE KOMPLEXEN GRÖSSEN IN IHRER BEZIEHUNG ZUR GEOMETRIE.}
\renewcommand{\thefootnote}{\arabic{footnote})}\setcounter{footnote}{0}
\begin{center}
\textbf{19.}\par\smallskip
\textls[100]{Weiterentwicklung der Lehre von den komplexen Grössen.}
\end{center}
\noindent\hspace*{\parindent}\textsc{Gauss} schliesst in der Selbstanzeige vom Jahre 1831 seine Auseinandersetzungen
über die imaginären Grössen mit den Worten: »Hier ist also die
Nachweisbarkeit einer anschaulichen Bedeutung von \(\surd{-1}\) vollkommen gerechtfertigt
und mehr bedarf es nicht, um diese Grösse in das Gebiet der
Gegenstände der Arithmetik zuzulassen« (W. II, S. 177). In ähnlicher Weise
hat er sich später (um 1850) in dem schon erwähnten Entwurf einer Abhandlung
über die Konvergenz der Reihen ausgesprochen: »Die imaginären Grössen
sind, solange ihre Grundlage immer nur in einer Fiktion bestand, in der
Mathematik nicht sowohl wie eingebürgert, als vielmehr nur wie geduldet
betrachtet, und weit davon entfernt geblieben, mit den reellen Grössen auf
gleiche Linie gestellt zu werden. Zu einer solchen Zurücksetzung ist aber
jetzt kein Grund mehr, nachdem die Metaphysik der imaginären Grössen in
ihr wahres Licht gesetzt und nachgewiesen ist, dass diese, ebenso gut wie die
negativen, ihre reale gegenständliche Bedeutung haben« (W. X\textsubscript{1}, S. 404).
\par
\textsc{Johann Bolyai} hat in einer 1837 verfassten, aber erst 1899 aus seinem
Nachlass veröffentlichten Schrift (Bol. II, S. 233) gegen die Ausführungen von
\textsc{Gauss} in der Selbstanzeige vom Jahre 1831 eine Reihe Einwendungen erhoben,
\par\smallskip\noindent{\scriptsize X\textsubscript{2} Abh. 4.\hfill 9}\par
\clearpage\SourceLeaf{397}\ArithmeticHead{66}{L}{STÄCKEL, GAUSS ALS GEOMETER.}
\noindent darunter auch die, dass \textsc{Gauss} »sich auf die Betrachtung des Raumes stütze,
die man in der Arithmetik vermeiden soll«, und er hatte selbst eine rein
arithmetische Einführung der komplexen Grössen gegeben, die im Wesentlichen
mit \textsc{Hamiltons}\footnote{R. W. \textsc{Hamilton}, \textit{Theory of conjugate functions, or algebraic couples}, Transactions of the Royal Irish Academy, Vol. 17, Dublin 1837, S. 393.} gleichzeitiger Begründung durch das Rechnen mit Grössenpaaren
übereinstimmt. Verschiedene Äusserungen von \textsc{Gauss} gestatten den
Schluss, dass auch er, eine 1831 im Keime vorhandene Auffassung weiter
entwickelnd, später zu einer von räumlichen Betrachtungen unabhängigen Auffassung
der komplexen Grössen übergegangen ist.
\par
In der Selbstanzeige vom Jahre 1831 wird ausgeführt, dass die komplexen
Grössen zur Darstellung der Relationen dienen können, die zwischen den Elementen
einer Mannigfaltigkeit von zwei Dimensionen stattfinden, und es heisst
dann, dass sich diese Verhältnisse nur durch eine Darstellung in der Ebene
zur Anschauung bringen liessen (W. II, S. 176). Noch entschiedener sagt
\textsc{Gauss} in dem zweiten Entwurf einer Abhandlung über die Konvergenz der
Reihen (um 1850): »Zuvörderst ist die bekannte Versinnlichung der komplexen
Grössen in Erinnerung zu bringen. . . . Es wird damit nur bezweckt, die
Bewegung in dem an sich vom Räumlichen unabhängigen Felde der abstrakten
komplexen Grössen zu erleichtern, und eine Sprache für dasselbe zu vermitteln«
(W. X\textsubscript{1}, S. 407).
\par
Diese Sprache für die Lehre von den »abstrakten« komplexen Grössen
hat \textsc{Gauss} in ihren Anfängen geformt. Eine nach der Stetigkeit fortschreitende
Reihe komplexer Grössen bildet einen \textls[100]{Zug}; jede der dem Zuge angehörigen
Grössen ist eine \textls[100]{Stelle} des Zuges. Ist der Zug geschlossen, so fügen sich
die nach der Stetigkeit zusammenhängenden komplexen Grössen, die in dem
Zuge ihre Begrenzung finden, zu einer \textls[100]{Schicht} zusammen. Man erkennt,
dass die geometrischen Namen Linie, Punkt, Fläche vermieden werden. In
einer Fussnote wird noch hervorgehoben, dass »die abstrakte allgemeine Lehre
von den komplexen Grössen mit der Wechselbeziehung zwischen vorwärts-rückwärts
und rechts-links nichts zu schaffen hat« (W. X\textsubscript{1}, S. 408).
\par
Was man vermisst, ist eine Erklärung, in welchem Sinne die formalen
Bildungen \(x+iy\) als \textls[100]{Grössen} bezeichnet werden dürfen. \textsc{Gauss} dürfte auch
hierüber seine Gedanken gehabt haben, denn in dem bereits angeführten Briefe
\clearpage\SourceLeaf{398a}\ArithmeticHead{67}{R}{DIE KOMPLEXEN GRÖSSEN IN IHRER BEZIEHUNG ZUR GEOMETRIE.}
\noindent an \textsc{Bessel} vom 21. November 1811, in dem er von den Funktionen einer
komplexen Veränderlichen spricht, sagt er: »Man sollte überhaupt nie vergessen,
dass die Funktionen, wie alle mathematischen Begriffszusammensetzungen,
nur unsere eigenen Geschöpfe sind und dass, wo die Definition,
von der man ausging, aufhört, einen Sinn zu haben, man eigentlich nicht
fragen soll: \textls[100]{Was ist?} sondern \textls[100]{was konveniert anzunehmen?} damit ich
immer konsequent bleiben kann. So z. B. das Produkt aus \(-\, .\, -\)« (W. X\textsubscript{1},
S. 363). Wenn man die Äusserungen über die allgemeine Arithmetik in der
Selbstanzeige vom Jahre 1831 hinzunimmt, wo das Gebiet der Zahlen stufenweise
erweitert wird (W. II, S. 175), so ergibt sich, wie nahe \textsc{Gauss} dem
Prinzip der Permanenz gekommen ist.
\clearpage\endgroup
